% language=us

\environment drawinglines-style

\startcomponent drawinglines-instances

\usemodule[scite]

\startchapter[title=Lua instances]

In various \LUAMETAFUN\ features we let \LUA\ do some calculations, in particular
math related ones, for instance when render functions or surfaces. That made me
wonder if using a dedicated \LUA\ instance would make sense, because there is
less interference of the garbage collector.

So I did a quick experiment and indeed there can be some gain but it is not
entirely trivial to integrate it in a transparent way, although most can be
hidden for the user. Given that the gain is not really putting a dent in
performance when we look at the whole run, for now it will not be applied. But we
can play a bit with the experimental feature as a math only sandbox.

Results can be piped back to the main instance, and one could preload functions
but a more extensive setup would also add some overhead, which then makes little
sense. Only the sandbox feature remains. As the experiment involves little code
we keep it for now. Here is the test:

\startbuffer
\startluacode
    lua.usestringcall(true)

    local n = 1 * 1000 * 1000
    local l = lua.luastringcall
    local m = lua.mpsstringcall

    local sin = math.sin
    local
    f = function()
        local a = sin(12)
        local b = sin(12)
        local c = a + b
     -- return 1, 2, 3, 4
        return c
     -- return c, true
    end

    local s = [[
        local sin = math.sin
     -- local
        f = function()
            local a = sin(12)
            local b = sin(12)
            local c = a + b
         -- return 1, 2, 3, 4
            return c
         -- return c, true
        end
    ]]

    -- register function (global)

    l(s)
    m(s)

 -- local s = s .. " " .. [[return f()]]

    local s = [[return f()]]
    local d = string.dump(load(s),true)

    s = d

    local c = os.clock() for i=1,n do local a = l(s) end print("l",l(s),os.clock()-c)
    local c = os.clock() for i=1,n do local a = m(s) end print("m",m(s),os.clock()-c)
    local c = os.clock() for i=1,n do local a = f(s) end print("f",f(s),os.clock()-c)
\stopluacode
\stopbuffer

\typebuffer[option=TEX]

On the console we see the following. Of course every call has to load the
(stripped) bytecode, which is why the direct call is so much more efficient.

\starttyping
l  -1.0731458360008699  0.952
m  -1.0731458360008699  0.889
f  -1.0731458360008699  0.159
\stoptyping

So, in order for this to work in our advantage we would have to register
functions in the \type {mps} namespace and communicate from the main \type {lua}
instance (where the user interface lives) that we'd like to call that specific
one. Because we use the same \LUA\ version for both the same bytecode can be used.

The only place where we trigger functions from \CCODE\ is in implicit three
dimensional surface generation. And complicating that even more that makes little
sense. Of course \typ {lua.luastringcall} might have some use of its own, who
knows. As a side note: the \type {return} in \type {f} in the example is quite
efficient because we pop return values from \type {mps} and push them to the main
instance.

\stopchapter

\stopcomponent

